\ifmostrarGabarito{
\begin{minted}{python}
# Resposta:

# 0 =< 15m3 , valor = 10
# > 15m3 e <= 40 , 1.2 por m3 consumido a mais
# >40m3, 3 por m3 consumido a mais

def calcular_conta_individual( consumo ):
    if (not type(consumo) == int) or (not type(consumo) == float):
        return -1
        
    if not (0 < consumo_1 < 5000):
        return -2

    if consumo > 40:
        excedente = consumo - 40
        return 10 + 25*1.2 + excedente*3
    elif 15 < consumo <= 40:
        excedente = consumo - 15
        return 10 + excedente * 1.2
    else: # taxa padrão da empresa
        return 10



def calcular_conta_das_familias( consumo_1 , consumo_2 , consumo_3 ):

    conta_residencia_1 = calcular_conta_individual(consumo_1)

    if conta_residencia_1 == -1:
        print("Argumento consumo 1 não é inteiro ou float")
        return -1

    if conta_residencia_1 == -2
        print(f"Consumo 1 ({consumo_1}) fora da faixa permitida.")
        return -2

    conta_residencia_2 = calcular_conta_individual(consumo_2)

    if conta_residencia_2 == -1:
        print("Argumento consumo 2 não é inteiro ou float")
        return -1

    if conta_residencia_2 == -2
        print(f"Consumo 2 ({consumo_2}) fora da faixa permitida.")
        return -2

    conta_residencia_3 = calcular_conta_individual(consumo_3)

    if conta_residencia_3 == -1:
        print("Argumento consumo 3 não é inteiro ou float")
        return -1

    if conta_residencia_3 == -2
        print(f"Consumo 3 ({consumo_3}) fora da faixa permitida.")
        return -2

    consumo_total = consumo_1 + consumo_2 + consumo_3
    valor_total = conta_residencia_1 + conta_residencia_2 + conta_residencia_3
    print(f"O valor consumido pela familia é de {consumo_total} m3.")
    print(f"Valor a ser pago é de {valor total} reais")
    return valor_total


\end{minted}
}